% ---------------------------------------------------------------------------
% appendix-tooling.tex
% ---------------------------------------------------------------------------
\chapter{Setting Up a Toolchain}
\label{app:tooling}

\chapabstract{What to install on your own machine so that nothing in this book
requires a licence, and what changes when you do have one.}

\section{The free toolchain}

Everything in \cref{ch:verilator}, \cref{ch:python} and most of
\refexamplebook runs on the following, all of which install from a package
manager on any recent Linux distribution and work under WSL2 on Windows.

\begin{center}
\small
\begin{tabular}{@{}llp{0.42\textwidth}@{}}
\toprule
\textbf{Tool} & \textbf{What it is} & \textbf{Why you want it}\\
\midrule
\tool{Verilator} & \SV{} to C++ compiler & fast simulation, and the best
  free linter for \SV{} by a wide margin\\
\tool{Icarus Verilog} & event-driven Verilog simulator & four-state,
  event-driven semantics closest to a commercial simulator; partial \SV{}\\
\tool{GTKWave} or \tool{Surfer} & waveform viewer & reads VCD, FST and GHW\\
\tool{cocotb}    & Python testbench framework & drives any of the above from
  Python\\
\tool{Python 3} with \tool{numpy} & golden models & \\
\tool{GNU Make}, \tool{g++} & build & Verilator needs a C++17 compiler\\
\bottomrule
\end{tabular}
\end{center}

\subsection{Debian, Ubuntu, WSL2}

\begin{shcode}
sudo apt update
sudo apt install -y verilator iverilog gtkwave \
                    build-essential python3 python3-pip python3-numpy
pip3 install --user cocotb cocotb-bus pytest
\end{shcode}

Check the versions. Distribution packages lag: \tool{Verilator} in particular
gained a great deal between 4.x and 5.x, and this book assumes 5.x throughout.

\begin{shcode}
verilator --version    # 5.0 or later
iverilog -V | head -1
cocotb-config --version
\end{shcode}

\begin{tipbox}{Build Verilator from source if the packaged version is old}
Verilator builds in a few minutes and is worth building. Clone the repository,
check out a release tag, then \code{autoconf \&\& ./configure \&\& make -j
\&\& sudo make install}. The 4.x series does not support \code{--binary},
\code{--timing}, classes or virtual interfaces, so half of
\cref{ch:verilator} will not work on it.
\end{tipbox}

\subsection{macOS}

\begin{shcode}
brew install verilator icarus-verilog gtkwave
pip3 install cocotb
\end{shcode}

\subsection{Windows}

Use WSL2 with an Ubuntu image and follow the Debian instructions. Native
Windows builds of these tools exist, but the surrounding flow -- make, shell
scripts, cocotb -- is much less painful under WSL.

\section{The commercial toolchain}

\subsection{QuestaSim / ModelSim}

Your department almost certainly has a licence, and Intel distributes a free
ModelSim edition with a size limit that is generous enough for anything in this
book. The one thing to check is that your version supports the \SV{} and UVM
features you need: UVM 1.2 support has been standard since roughly 10.x, and
IEEE 1800.2 since 2020.x.

\begin{shcode}
which vlog vsim vlib
vlog -version
\end{shcode}

\subsection{Getting UVM}

Three options, in decreasing order of convenience.

\begin{enumerate}
  \item Use the simulator's built-in copy: \code{vlog -sv +define+UVM\_NO\_DPI}
        is not needed if you compile with \code{-uvm} or set
        \code{-uvmhome} to the vendor's bundled tree.
  \item Download the Accellera reference implementation, set
        \code{UVM\_HOME}, and compile \file{\$UVM\_HOME/src/uvm\_pkg.sv} with
        \code{+incdir+\$UVM\_HOME/src}. Add the DPI library if you want
        regular expression matching in the config database.
  \item \tool{pyuvm}, if you would rather write UVM in Python. See
        \cref{ch:python}.
\end{enumerate}

\section{A project layout that works}

The layout used by the accompanying sources, and a reasonable default for a
thesis or a lab project. It is organised by \emph{language}, because the same
design exists here in two of them, and every chapter of this book that names a
file names it with one of these prefixes.

\begin{txtcode}
code/cordic/
  sv/                     synthesisable SystemVerilog, the interface,
                          the class testbench
    cordic_pkg.sv  cordic_rot.sv  cordic_if.sv
    tb_cordic_pkg.sv  tb_cordic_top.sv
    cordic_atan_rom.svh  cordic_vectors.svh   (generated)
  vhdl/                   synthesisable VHDL and the VHDL testbench
    cordic_pkg.vhd  cordic_rot.vhd  tb_cordic_rot.vhd
  verilator/              C++ testbenches and the svtlm.h mini-framework
    tb_cordic.cpp  tb_vectors.cpp  main_uvmlike.cpp
    cordic_env.h  svtlm.h  cordic_vectors.h     (generated)
  python/                 golden model, vector generator, cocotb tests
    cordic_model.py  gen_vectors.py  test_cordic.py
    Makefile.cocotb
  vectors/                generated, safe to delete
    cordic_vectors.txt
  build/                  build outputs, not committed
  Makefile
\end{txtcode}

Two things are worth copying. Generated files sit next to the files that
consume them and carry a \code{GENERATED} banner in their first line, so nobody
edits one by hand. Everything under \file{vectors/} and \file{build/} can be
deleted at any moment and rebuilt with \code{make vectors}.

If you prefer to split by role rather than by language, the usual alternative
is \file{rtl/} for synthesisable sources and \file{tb/} for everything else.
Either works. What does not work is mixing them, so decide on the first day.

\section{Continuous integration}

The single most useful thing you can add to a student project, and it takes
ten lines.

\begin{txtcode}
name: rtl
on: [push, pull_request]
jobs:
  check:
    runs-on: ubuntu-latest
    steps:
      - uses: actions/checkout@v4
      - run: sudo apt-get update && sudo apt-get install -y verilator iverilog
      - run: pip install cocotb
      - run: make lint
      - run: make sim
\end{txtcode}

The \code{lint} target should be \code{verilator --lint-only -Wall} over every
synthesisable file, and it should fail the build. Treat a lint warning the way
you treat a VHDL analysis error, because in this language that is what it is.

\section{Editor support}

\tool{Verible} provides a formatter and a language server for \SV{}; \tool{rust\_hdl}
does the same for VHDL. Both integrate with any editor that speaks LSP. A
formatter matters more in \SV{} than in VHDL because the language gives you
so much more freedom to lay code out badly.
